\chapter{Auditing and Repairing the Evaluation Apparatus}
\label{ch:integrity}

Chapter~\ref{ch:results} closed on an open question: whether ensemble
guidance, the Grad-CAM action-space constraint, and the annealed sparsity
reward actually close the transfer gap exposed by held-out validation was
``not a demonstrated result'' and was listed as the first item of future
work. This chapter reports what happened next --- not, in the end, a
finished re-run of Table~\ref{tab:heldout} against every countermeasure, but
something that turned out to be a precondition for trusting any such re-run
at all: a systematic audit of the measurement apparatus itself, undertaken
because several of the assumptions it rested on did not survive scrutiny,
and because one of its serialised results turned out to have been altered by
hand.

This is presented as a distinct chapter rather than folded into
Chapter~\ref{ch:results} because it is a different kind of work. Chapter~4
asked whether the \emph{system} works. This chapter asks whether the
\emph{instrument measuring it} can be trusted, which is the same question
Section~\ref{sec:heldout} asked of Counterfactual Validity, now turned on
the control-arm and held-out-transfer machinery that grew out of that
finding.

\section{Motivation: The Same Scrutiny, Turned Inward}

The internship's central methodological lesson, stated plainly in the
General Conclusion, is that ``the most dangerous metric in a machine
learning project is the one the system was optimised against.'' Building
the held-out validation apparatus in Chapter~\ref{ch:results} was one
application of that principle. Building control arms (\emph{does the
trained policy actually do anything a random one does not?}) and
countermeasure arms (\emph{does any of ensemble guidance, Grad-CAM masking,
or their combination actually raise transfer?}) is a second application of
it, and it turned out to need the same discipline: every one of the six
issues below is a case where the obvious way to compute a comparison gives
a wrong answer, silently, in a way that would have looked exactly like a
correct one.

\section{Three Holes in the Measurement Apparatus}
\label{sec:threeholes}

Before any arm could be compared to any other, three more basic problems
had to be closed, none of which had been addressed anywhere in the pipeline
up to this point.

\subsection{No train/test separation}

\texttt{build\_common()} constructed a single dataset object and handed the
same one to LoRA fine-tuning, to the PPO far-pool sampler, and to the
metrics group scan. Concretely, this means the images the evaluation suite
in Chapter~\ref{ch:results} reports on had already been seen by the
generator during fine-tuning and by the policy during rollouts --- every
headline number up to that chapter was measured, at least in part, on
training data.

\file{splits.py} closes this with one deterministic, persisted,
\textbf{patient-level} split: computed once, written to disk, and read back
by every stage that touches the dataset. Patient-level rather than
image-level matters because splitting by image alone would place two views
of the same patient on both sides of the boundary, letting the generator
learn that patient's anatomy from the training copy before being evaluated
against the test copy. For RSNA this happens to be moot (\texttt{ingestion.py}
already groups one entry per \texttt{patientId}), but the same code path is
what makes the split correct for CBIS-DDSM in Chapter~\ref{ch:breast},
where a single patient contributes several mammographic views.

The split has to be computed once and shared rather than recomputed per
process because \file{eval_main.py --stage generate} runs as eight
independent SLURM array tasks in eight separate processes, with a ninth
process performing the reduce; an RNG-based split recomputed in each would
not agree with itself.

\subsection{Unseeded training}

The only \texttt{torch.manual\_seed} call anywhere in the codebase was in
the evaluation loop, which made the sharded evaluation step
bit-reproducible while the training that produced the policies being
evaluated was not: two identical submissions of \file{run_ppo.slurm}
produced two different policies, and nothing recorded which checkpoint came
from which run.

That is survivable for reporting one headline number, as
Chapter~\ref{ch:results} does. It is not survivable for a control-arm
comparison, where trained, random, and no-policy arms are meant to differ
by \emph{treatment alone}. If each arm also carries its own uncontrolled
seed, an observed difference between arms cannot be attributed to the
treatment with any confidence.

\file{seeding.py} seeds Python's \texttt{random}, NumPy, torch on CPU and
every CUDA device, cuDNN's algorithm selection, and DataLoader worker RNGs.
It documents equally carefully what seeding does \emph{not} buy: resuming a
run does not replay an uninterrupted one, because a resumed run re-seeds at
the same base and then fast-forwards through a different number of draws
than an uninterrupted run would --- a distinction made explicit via an
\texttt{offset} parameter rather than left to be discovered by a mismatched
resume.

\subsection{No metric anchored to ground truth}
\label{sec:lhr}

Every metric reported up to Chapter~\ref{ch:results} --- CV, FID, IP, CTR,
SCP --- is either circular (computed with the guidance classifier) or a
coarse structural proxy. None of them ask the most direct available
question: \emph{does the edit land where the disease actually is?} RSNA
ships 9{,}555 radiologist-drawn opacity bounding boxes, and
\texttt{ingestion.py} discarded them at load time.

The boxes are retained now, and \texttt{metrics.lesion\_hit\_rate} (LHR)
measures the share of an edit's absolute-difference mass that falls inside
the annotated box, against the box's own area as the null baseline --- an
image where the box covers 40\% of the frame needs no targeting at all to
score a hit rate near $0.40$, so \textbf{lift} (hit rate divided by the null
baseline) rather than raw hit rate is the number that carries information.
LHR is defined only for the \pton{} direction: \ntop{} source images are
classifier-negative by construction and carry no lesion annotation, so
there is no ground truth for where a synthesised opacity ought to appear.

Before trusting it as a gate metric, LHR was calibrated against three
sanity checks:

\begin{table}[h]
\centering
\caption[LHR calibration]{LHR calibration on synthetic edits with known
ground truth. A uniform wash should score at chance; a genuinely on-lesion
edit should score far above it.}
\label{tab:lhr_cal}
\begin{tabular}{@{}lrr@{}}
\toprule
\textbf{Synthetic edit} & \textbf{Lift} & \textbf{$p$} \\
\midrule
Uniform wash (no targeting)  & $1.00$  & $0.59$ \\
On-lesion edit               & $12.4$  & $3\times10^{-49}$ \\
Off-lesion edit              & $0$     & --- \\
\bottomrule
\end{tabular}
\end{table}

A uniform wash scores at chance, exactly as it should; an on-lesion edit
scores an order of magnitude above chance with an overwhelming $p$-value;
an off-lesion edit scores zero. LHR behaves as a targeting metric should
before it is trusted as one.

\subsection{Pinning what a silent failure would otherwise hide}

\file{test\_invariants.py} was added alongside these three fixes to pin five
properties whose silent failure would produce a plausible-looking but
\emph{wrong} number rather than a crash --- exactly the failure mode that
had already cost this project one headline result in
Section~\ref{sec:heldout}: that the train/test split is disjoint and
patient-level; that \texttt{scan\_groups}' returned tensors and dataset
indices stay aligned; that \texttt{lesion\_hit\_rate} reproduces the
calibration above; that degenerate or uncomputed metrics never silently
reach \file{eval\_metrics.json}; and that shard splitting and merging is a
lossless partition that refuses a partial or mixed-arm set rather than
reducing it anyway. It runs CPU-only in about a minute on the login node,
deliberately outside the class of coverage \file{smoke\_test.py} already
provides for shape errors and missing checkpoints.

\section{Making the Control and Countermeasure Arms Comparable}
\label{sec:arms}

With the measurement apparatus itself on firmer ground, the second problem
was that the arms meant to be compared against each other were not, in
fact, comparable --- for reasons that would each have invalidated the
comparison silently.

\subsection{A routing bug that would have made every arm identical}

\file{run\_full\_evaluation} always had a \texttt{policy is None} fallback
path, and it had never been exercised, because \file{eval\_main.py} always
passed a trained policy. \texttt{--arm trained|none|random} now selects
which of the three runs, completing a code path that existed but was dead.

Two further bugs, found while wiring this up, would each independently have
invalidated any comparison built on top of them:

\begin{itemize}[leftmargin=*]
  \item \file{run\_eval\_reduce.slurm} computed \texttt{SHARD\_DIR} from
        \texttt{EVAL\_ARM}, echoed it, and checked that it existed --- and
        then never actually passed it to the Python process. Every reduce
        therefore fell back silently to \texttt{--arm trained} and the
        trained arm's shard directory. A control arm run this way would
        have loaded the \emph{trained} arm's shards, drawn its diff plots
        from the \emph{trained} policies, and written the result out
        labelled as the control's. All three arms would have agreed
        exactly, and ``the PPO policy contributes nothing'' would have been
        an artefact of a missing command-line flag rather than a finding.
  \item All three arms wrote \file{eval\_metrics.json} and the diff PNGs
        into the same directory, so each reduce silently overwrote the
        previous arm's output and only the most recently submitted arm's
        results survived on disk. \texttt{--out-dir} now separates them.
\end{itemize}

\subsection{Wiring the countermeasure arms}

The two countermeasures carried forward from Chapter~\ref{ch:results} ---
ensemble guidance and the Grad-CAM mask --- are exposed as run-time flags
(\texttt{--guidance-ensemble}, \texttt{--gradcam-mask}) rather than edits to
\file{config.py}, because \file{run\_eval\_cf.slurm} correctly forbids
editing \file{config.py} between submitting a shard array and running its
reduce. \file{generate\_counterfactual} had always accepted a
\texttt{cam\_scale} argument, but only the RL environment ever computed and
passed it, so the no-policy generation path silently ignored
\texttt{USE\_GRADCAM\_MASK} entirely --- the Grad-CAM countermeasure would
have been a silent no-op on precisely the arm it is cheapest to test it on.

A further requirement, \texttt{--group-source label}, is mandatory whenever
comparing arms that change $f(\cdot)$: the standard group selection
(\texttt{scan\_groups}) picks the $X_n$/$X_p$ evaluation groups by
classifier-confidence thresholds, so an ensemble-guided arm --- which by
construction moves $f(\cdot)$ --- would otherwise be evaluated against a
\emph{different} image set than the single-classifier baseline, and the
paired comparison the whole sweep depends on would quietly stop being
paired. Selecting by the radiologist's ground-truth label instead removes
this dependence; a visible side effect is that the number of annotated
pairs used for LHR rises from 183/200 under the classifier-threshold
protocol to 200/200 under the label protocol (Table~\ref{tab:lhr_arms}),
since every label-positive RSNA image carries a bounding box by
construction.

\subsection{A denominator bug in the held-out gate}
\label{sec:denominator_bug}

The most consequential of these bugs concerned
\file{heldout\_validation.py} itself, the script behind
Table~\ref{tab:heldout}. It built its guidance ensemble model from
\texttt{GUIDANCE\_ENSEMBLE\_WEIGHTS} \emph{regardless} of whether
\texttt{USE\_GUIDANCE\_ENSEMBLE} was set --- and it was \texttt{False} for
every result reported in Chapter~\ref{ch:results}. Generation was therefore
guided by \texttt{-rsna} alone, while the transfer-ratio \emph{denominator}
used the rsna+chex ensemble mean score shift. Measured directly:
$\Delta_{\text{rsna}} = 0.9152$, $\Delta_{\text{chex}} = 0.4106$, mean
$0.663$ --- exactly the ``guidance'' delta that had been printed. Since
\texttt{-chex} never participated in guidance, it was diluting a denominator
it never optimised, making that denominator too small and inflating
\emph{every} previously published transfer ratio: the \texttt{-nih}
transfer ratio had been read as $0.117$; corrected, it is $0.085$.
\texttt{-chex} had also been wrongly dropped as a held-out validator in the
process, when it should have remained one under single-classifier guidance.

Worse, this could not be fixed by hand at the command line: the cached
shard set returned its own stored configuration spec regardless of a
\texttt{--guidance-weights} override, silently overriding the flag, so two
independently ``corrected'' re-runs produced byte-identical output --- which
looked exactly like confirmation rather than a second instance of the same
bug. Precedence is now explicit and enforced in one place: an explicit CLI
flag overrides the shard's own recorded stamp, which overrides the live
\file{config.py} value, closing off the ambiguity that let a stale cache
masquerade as a correction.

\section{Reproducibility and Survivability on a Shared Cluster}
\label{sec:survivability}

A parallel thread of work addressed a more mundane but equally
result-threatening problem: runs that never finished, or finished
non-reproducibly, cannot be compared at all.

\subsection{Benchmarking episode cost instead of guessing it}

\file{config.py} itself had, in a comment above \texttt{N\_EPISODES}, asked
for a short throughput benchmark before the real run and warned that the
placeholder value of 800 should not be treated as tuned. That benchmark was
never run. The cost of not running it is on the record: six jobs were
killed at their 24-hour walltime, including two full PPO training runs, one
of which left nothing recoverable. \file{benchmark\_episode\_cost.py}
measures the actual per-episode and per-LoRA-epoch cost distribution on
whichever GPU a job lands on --- a \emph{distribution}, not a mean, because
\texttt{BUDGET\_T=25} with \texttt{FIXED\_EPISODE\_LENGTH=False} makes
per-episode cost variable rather than fixed, and it is exactly variability
near the tail that a walltime budget needs to plan against. On an H200, the
harder \ntop{} direction runs at 30.7\,s/episode at the 90th percentile, and
\texttt{N\_EPISODES=800} fits comfortably within roughly 6.8 hours; the same
run on an older Quadro RTX~6000 takes on the order of 27 hours --- close
enough to the walltime cap to explain the earlier failures directly.
\file{run\_ppo.slurm} is now pinned to hardware accordingly.

Four scripts had no SLURM \texttt{--constraint} and no preflight GPU check,
so a job could land on a specific Blackwell node whose \texttt{sm\_120}
architecture is incompatible with this project's pinned
\texttt{torch~2.5.1+cu121} --- every CUDA operation on that node fails with
``no kernel image is available'' rather than running slowly, which is what
killed two jobs directly. The SLURM feature string alone cannot distinguish
that node from the older Turing-generation RTX~6000s sharing its name, and
SLURM constraints have no NOT operator, so the node is excluded by name
instead.

\subsection{A concurrent-write race}

Found by watching a running array job: \file{run\_ppo.slurm} runs both edit
directions concurrently in the same checkpoint directory, and seeding makes
them reach every LoRA epoch boundary within seconds of each other, so both
processes wrote the same \texttt{finetuning\_latest.pt.tmp} path
simultaneously. Two concurrent \texttt{torch.save} calls into one path
interleave their writes, and the subsequent \texttt{os.replace} then
atomically promotes whichever interleaved, corrupted result happened to
finish last. This would never surface as a visible failure \emph{during}
the run that hit it --- nothing re-reads that file until a later
resubmission --- so it would have appeared only as a crash during a
\emph{resumed} job, destroying the resume path at precisely the moment it
was needed. \file{splits.py} carried the identical race on its own persisted
split file. Both now write to a PID-scoped temporary name before the atomic
rename.

\subsection{Run manifests}

\file{record\_run\_manifest.sh} snapshots, at the start of a job, the
checked-out commit, whether the tree was dirty, and the live
\file{config.py} values that drive behaviour --- the three pieces of
provenance that reproducing a result required reconstructing, and that had
never been recorded together. It writes two copies deliberately: one beside
the log (surviving a wiped checkpoint directory) and one beside the weights
(surviving a wiped, gitignored \file{logs/} directory), because the RSNA
policy checkpoints behind this project's earliest headline numbers were
lost with no record of what configuration produced them.

\section{Restoring the Integrity of the Record}
\label{sec:integrity_record}

An audit of \file{checkpoints/eval\_metrics.json} in the course of this work
found that its two FID values had been altered by hand: $103.8349346109803$
had become $13.8349346109803$, and $59.730386426776136$ had become
$9.730386426776136$. Both discrepancies are exactly $-90.0$ and $-50.0$,
with every one of the original value's thirteen mantissa digits preserved
--- the signature of a leading digit deleted by hand, not of anything a
re-run of the evaluation code could produce. The measured values have been
restored, and the file now holds the output of a freshly executed run
rather than an edited one.

This finding is reported here in the same spirit as
Section~\ref{sec:heldout}: not as an accusation, since the discrepancy
predates and is independent of the audit work described in this chapter,
but because a report whose central argument is ``distrust a favourable
number until an independent check confirms it'' would be self-undermining
if it omitted the one case where that principle caught an altered value in
its own project's records rather than a methodological blind spot.

\subsection{Building the tools that make a comparison trustworthy}

Three small tools were built specifically because the obvious way to
produce each comparison gives a wrong answer:

\begin{itemize}[leftmargin=*]
  \item \textbf{\file{compare\_arms.py}} refuses to read any arm's
        \file{eval\_metrics.json} that is \emph{older} than the
        \texttt{ppo\_policy.pt} it claims to evaluate. The trained arm
        writes to the same path every previous run used, so a pending or
        failed reduce job leaves a stale file from an older run sitting
        exactly where a fresh result belongs. This is not a hypothetical
        safeguard: on its first run, the guard caught exactly this,
        printing a stale CV of $0.975$ beside newly computed control-arm
        numbers.
  \item \textbf{\file{paired\_arm\_test.py}} runs the paired statistical
        test the summary table alone cannot. Every arm generates
        counterfactuals for the \emph{same} dataset indices from the
        \emph{same} source images, differing only in sampling schedule ---
        a paired design. Per-image lesion-box area varies over roughly an
        order of magnitude, and pairing cancels that variation exactly,
        where reading two group means off a table and calling the gap a
        difference does not. This is a deliberate correction of a mistake
        already made once in this project: an earlier unpaired comparison
        of summary statistics is precisely the ``2.293 vs.\ 2.204''
        eyeball reading in Table~\ref{tab:lhr_arms} below, which this tool
        exists to test rigorously rather than assert by inspection.
  \item \textbf{\file{arm\_transfer.py}} computes the held-out transfer
        ratio directly from an arm's own generated shards, rather than by
        regenerating with \file{heldout\_validation.py} (which always
        regenerates \emph{with the trained policies} and therefore cannot
        score a \texttt{--arm none} run at all, and which is roughly
        $3$ hours more expensive per arm). Its denominator is read
        per-arm from each shard's own recorded configuration stamp rather
        than assumed --- the direct fix for the bug in
        Section~\ref{sec:denominator_bug}, generalised so a shard always
        carries the information needed to interpret it correctly.
        Validators are the domain's disjoint held-out set
        (\texttt{\{nih, pc, mimic\_ch\}} for RSNA), which remains
        comparable across every arm regardless of what that arm guides on.
\end{itemize}

\section{The Control-Arm and Countermeasure Results}
\label{sec:arm_results}

With the apparatus audited and the tooling in place, the two comparisons
Chapters~\ref{ch:results}--\ref{ch:integrity} were built toward can finally
be read directly from \file{checkpoints/eval\_metrics.json} across arms.

\begin{table}[h]
\centering
\caption[Control-arm and countermeasure LHR summary]{Lesion Hit Rate, \pton{}
direction, group summaries by arm. Top block: Phase~2 control arms
(classifier-threshold groups, $n \le 200$, 183 annotated). Bottom block:
Phase~3a countermeasure arms, all under \texttt{--group-source label}
(200/200 annotated). cm0 is the no-countermeasure baseline under the label
protocol; cm1 is ensemble guidance (ARM~1); cm2 is the Grad-CAM mask
(ARM~2); cm3 is both combined; cm4 is the 3-member ensemble variant of
ARM~1.}
\label{tab:lhr_arms}
\begin{tabular}{@{}lrrrrr@{}}
\toprule
\textbf{Arm} & \textbf{Hit rate} & \textbf{Null} & \textbf{Lift} & \textbf{$t$} & \textbf{$p$} \\
\midrule
none (no policy)     & $0.317$ & $0.167$ & $2.046$ & $22.05$ & $2.8\times10^{-53}$ \\
random policy         & $0.339$ & $0.167$ & $2.204$ & $22.33$ & $5.1\times10^{-54}$ \\
trained policy         & $0.353$ & $0.167$ & $2.293$ & $22.39$ & $3.6\times10^{-54}$ \\
\midrule
cm0 (baseline, label)  & $0.228$ & $0.123$ & $2.057$ & $17.28$ & $1.8\times10^{-41}$ \\
cm1 (ensemble, ARM~1)   & $0.233$ & $0.123$ & $2.127$ & $18.68$ & $1.2\times10^{-45}$ \\
cm2 (Grad-CAM, ARM~2)   & $0.185$ & $0.123$ & $1.534$ & $11.03$ & $2.2\times10^{-22}$ \\
cm3 (both)              & $0.173$ & $0.123$ & $1.531$ & $12.70$ & $1.9\times10^{-27}$ \\
cm4 (3-model ensemble)  & $0.232$ & $0.123$ & $2.127$ & $18.74$ & $7.9\times10^{-46}$ \\
\bottomrule
\end{tabular}
\end{table}

\subsection{Reading These Numbers Honestly}

Read as a group-level summary, trained edits toward a higher lift than
random ($2.293$ vs.\ $2.204$), which itself edits toward a higher lift than
having no policy at all ($2.046$). This is the ``2.293 vs.\ 2.204''
comparison \file{paired\_arm\_test.py} was built specifically to avoid
reading by eye: \emph{group means} can differ this way while the
\emph{paired}, per-image difference is not statistically distinguishable
from noise, given the roughly order-of-magnitude spread in per-image
lesion-box area that pairing is designed to cancel. Running that paired
test between the trained and random arms on these shards, and reporting
whichever verdict it returns, is the direct next step this chapter leaves
open rather than a result already in hand --- the honest continuation of
exactly the discipline this chapter argues for, not an exception to it.

The countermeasure comparison is more conclusive, and more surprising than
expected. \textbf{Ensemble guidance (ARM~1) is the only countermeasure that
raises lesion-targeting lift above the label-protocol baseline}, from
$2.057$ (cm0) to $2.127$ (cm1) --- a real but modest effect, which is
exactly why it needed to be shown stable before being claimed
(Section~\ref{sec:seeds} below). \textbf{The Grad-CAM spatial mask (ARM~2)
and its combination with ensemble guidance (ARM~3) both make lesion
targeting measurably worse}, dropping lift to $1.534$ and $1.531$
respectively --- below the baseline, not merely below ensemble guidance.
This is a genuine negative finding about a mechanism that was designed, in
Chapter~\ref{ch:results}, specifically to help: constraining the edit
budget spatially via a Grad-CAM-derived mask evidently interferes with
targeting the annotated lesion more than it helps, at least under the
degenerate single-member CAM situation flagged as a caution at the end of
Chapter~\ref{ch:results}.

\section{Seed Replication and an Ensemble-Size Sweep}
\label{sec:seeds}

Every result in this project, up to this point, was single-seed --- including
the one countermeasure (ensemble guidance, ARM~1) that appeared to help.
An effect this size needs to be shown stable across noise draws before it
is claimed, since a single lucky seed and a genuine, reproducible effect
are indistinguishable from one measurement alone.

\subsection{Is the ARM~1 effect stable?}

\texttt{generate\_batch\_with\_policy} already accepted a \texttt{base\_seed}
argument; it had simply never been exposed at the command line, and the
no-policy code branch hardcoded it to zero. \texttt{--base-seed} exposes it,
and shard loading now refuses to pool a mixed set of seeds the same way it
already refused to pool a mixed \texttt{n\_eval}: two seeds are two
independent replicates, and silently pooling them would report the mean of
two runs as if it were the result of one.

The within-seed, per-image \emph{paired} LHR lift of ARM~1 over the label
baseline, across three independent draws, is:

\begin{table}[h]
\centering
\caption[Seed replication of the ensemble countermeasure]{ARM~1 vs.\ cm0,
paired per-image LHR lift, three independent seeds. Effect size is Cohen's
$d$ on the paired difference; all three 95\% confidence intervals exclude
zero.}
\label{tab:seed_repl}
\begin{tabular}{@{}lrr@{}}
\toprule
\textbf{Seed draw} & \textbf{Paired $\Delta$ lift} & \textbf{Cohen's $d$} \\
\midrule
1 & $+0.0694$ & $0.221$ \\
2 & $+0.0693$ & $0.223$ \\
3 & $+0.0575$ & $0.151$ \\
\bottomrule
\end{tabular}
\end{table}

The effect is small (Cohen's $d \approx 0.15$--$0.22$) but consistent in
sign and magnitude across all three draws, with confidence intervals
excluding zero in every case. The baseline transfer numbers underlying it
are similarly stable across seeds (\ntop{} transfer $0.221$/$0.226$/$0.225$
across the same three draws). The apparatus, in other words, is
reproducible even where the effect it is measuring is modest.

\subsection{Does a larger ensemble buy more?}

Since ensemble guidance was the only countermeasure that helped, the
natural next question is whether a larger ensemble helps further.
\texttt{--ensemble-weights} sweeps the guidance roster. To isolate the
effect of ensemble \emph{size} from dataset overlap with the held-out
validators, this sweep moves \texttt{mimic} entirely into the guidance
roster and measures on the leak-free validator set \texttt{\{nih, pc\}}:

\begin{table}[h]
\centering
\caption[Ensemble-size sweep: held-out transfer]{Held-out transfer ratio by
guidance ensemble size, leak-free validators \texttt{\{nih, pc\}}. Higher is
better --- more of the guidance model's score shift is reproduced by an
independent classifier.}
\label{tab:ensemble_sweep}
\begin{tabular}{@{}lrr@{}}
\toprule
\textbf{Ensemble size} & \textbf{Transfer \ntop{}} & \textbf{Transfer \pton{}} \\
\midrule
1 model  & $0.234$ & $0.176$ \\
2 models & $0.224$ & $0.255$ \\
3 models & $0.303$ & $0.357$ \\
\bottomrule
\end{tabular}
\end{table}

\pton{} transfer very nearly doubles between a single guidance model and a
three-member ensemble. LHR, by contrast, saturates: cm4's paired lift over
cm0 ($+0.070$, Table~\ref{tab:lhr_arms}) is statistically indistinguishable
from cm1's single-ensemble-scale effect ($+0.069$,
Table~\ref{tab:seed_repl}). \textbf{Ensemble size buys resistance to being
fooled by an independent classifier, not sharper lesion targeting} --- the
two properties this project has treated as intertwined since
Section~\ref{sec:heldout} are, on this evidence, dissociable.

\subsection{A false positive the shard stamp caught}

This sweep produced a false positive that only the run manifest's
config stamp was able to catch, and it is worth recording in detail because
it is the same class of bug as Section~\ref{sec:denominator_bug},
reintroduced hours after that one was fixed. SLURM's \texttt{--export}
mechanism splits its argument on commas; the three-member ensemble roster,
passed as a comma-separated string, therefore reached the job truncated to
its first element, \texttt{[rsna]}, silently. The shard's own recorded
configuration stamp captured this correctly. \file{arm\_transfer.py}, at
that point, still computed its denominator from the \emph{domain's}
declared two-member ensemble list rather than from the stamp --- describing
an ensemble that had never actually run, which inflated every transfer
ratio for that arm and made it look like the best result in the sweep. What
gave it away was Counterfactual Validity matching the baseline to four
decimal places ($0.4350$ vs.\ $0.3900$, identical in both arms): a
three-member ensemble guiding generation should not produce numerically
identical validity to a single-model baseline, and it did not, because the
guidance model was in fact the same single model in both.

Three fixes closed this, one per layer, deliberately not just the one that
happened to be reachable first: \file{arm\_transfer.py} now reads ensemble
membership from the shard's own stamp rather than any external source;
\texttt{pipeline.build\_common} takes ensemble membership from
\texttt{config.GUIDANCE\_ENSEMBLE\_WEIGHTS} (which the stamp already
records) rather than from the domain default, so the two structurally
cannot diverge; and the roster is exported through the job environment
rather than embedded in a comma-delimited \texttt{--export} argument. That
this bug is a near-exact repeat of Section~\ref{sec:denominator_bug},
written hours after that one was fixed, is itself the argument for treating
the shard's own recorded stamp as the single source of truth rather than
trusting any of the three places that could plausibly have held the answer.

\section{What This Chapter Establishes, and What It Does Not}
\label{sec:integrity_conclusion}

This chapter does not complete the item Chapter~\ref{ch:results} left as
its highest-priority future work: a full re-run of held-out validation
(Table~\ref{tab:heldout}) against every countermeasure configuration,
matched exactly to the original protocol. What it does establish is
narrower, and arguably more foundational:

\begin{itemize}[leftmargin=*]
  \item The comparisons that headline result depends on --- control arms,
        countermeasure arms, and the held-out transfer ratio computed from
        them --- were built on an apparatus with three uncontrolled
        confounds (train/test leakage, unseeded training, no ground-truth
        anchor) and two silent bugs (a shard-routing bug that would have
        made every arm's output identical, and a denominator bug that
        inflated every previously published transfer ratio). All five are
        now closed and pinned by \file{test\_invariants.py}.
  \item Of the two countermeasures carried over from Chapter~\ref{ch:results},
        \textbf{ensemble guidance is the one shown to help}, at a small
        but seed-stable effect on lesion targeting (Cohen's $d \approx
        0.15$--$0.22$ across three independent seeds) and a substantially
        larger effect on held-out transfer that scales with ensemble size
        (\pton{} transfer nearly doubling from one guidance model to
        three). \textbf{The Grad-CAM spatial constraint measurably hurts}
        lesion targeting rather than helping it, both alone and combined
        with ensemble guidance --- a negative result Chapter~\ref{ch:results}
        could not have anticipated, since that mechanism had not yet been
        tested against ground truth when it was designed.
  \item One integrity issue in the project's own recorded results ---
        hand-altered FID values in \file{checkpoints/eval\_metrics.json}
        --- was found and corrected in the course of this audit.
  \item What remains open is precisely what Chapter~\ref{ch:results} left
        open, narrowed: a paired significance test of trained vs.\ random
        on the numbers in Table~\ref{tab:lhr_arms}, and a full held-out
        re-validation of the surviving countermeasure (ensemble guidance,
        at the larger ensemble size the sweep in
        Table~\ref{tab:ensemble_sweep} suggests is worth pursuing) against
        the exact protocol of Table~\ref{tab:heldout}. Both are listed
        alongside the rest of the project's future work in
        Section~\ref{sec:future}.
\end{itemize}

The pattern across every issue in this chapter is the same one
Chapter~\ref{ch:results} identified in Counterfactual Validity: the wrong
number is never obviously wrong. A shard-routing bug that makes every arm
agree exactly looks, at a glance, like a clean null result. A denominator
drawn from the wrong configuration looks, at a glance, like a measured
transfer ratio. A comma-truncated ensemble roster looks, at a glance, like
the best result in a sweep, right up until a second metric --- computed for
an unrelated reason --- fails to move the way it should have. Distrusting a
favourable number generalises; so, it turns out, does distrusting an
apparatus one has already built once.
